\begin{table}[t]
\centering
\renewcommand{\arraystretch}{1.25}
\begin{tabularx}{\linewidth}{l c X}
\hline
\textbf{Term} & \textbf{Symbol} & \textbf{Definition / Project Context} \\
\hline

Obukhov length
 & $L$
 & Scale at which shear and buoyancy production balance; fundamental length scale for SBL stability. \\

Dimensionless height
 & $\zeta$
 & $\zeta = z/L$. Non-dimensional height used in MOST functions; $\zeta>0$ in stable conditions. \\

Geometric mean height
 & $z_g$
 & $z_g = \sqrt{z_0 z_1}$. Representative height for log-layer profiles and where $\mathrm{Ri}_g$ is often evaluated. \\

Logarithmic mean height
 & $z_L$
 & $z_L = (z_1 - z_0) / \ln(z_1 / z_0)$. Exact representative height for the velocity difference in a logarithmic wind law. \\

Gradient Richardson number
 & $\mathrm{Ri}_g$
 & $\displaystyle \mathrm{Ri}_g(z)=\frac{(g/\theta)\,(\partial\theta/\partial z)}{(\partial U/\partial z)^2}$. Local stability diagnostic used to determine the collapse threshold $\mathrm{Ri}_c$. \\

Bulk Richardson number
 & $\mathrm{Ri}_b$
 & $\displaystyle \mathrm{Ri}_b=\frac{(g/\bar{\theta})\,(\Delta\theta\,\Delta z)}{(\Delta U)^2}$. Layer-averaged stability measure and main input for NWP closures. \\

Bias ratio
 & $B$
 & $B = \mathrm{Ri}_g(z_g) / \mathrm{Ri}_b$. Diagnostic of curvature-induced bias. $B>1$ indicates that $\mathrm{Ri}_b$ underestimates local stability. \\

Grid damping factor
 & $G$
 & $G(\zeta,\Delta z)$; multiplicative correction applied to eddy diffusivity: $K_{\text{new}} = K_{\text{old}} \cdot G$. ML model predicts $G$ from $(\zeta,\Delta z)$. \\

Eddy diffusivities
 & $K_m, K_h$
 & Momentum and heat diffusivities (m$^2$ s$^{-1}$). Grid damping factor $G$ is applied directly. \\

Critical Richardson number
 & $\mathrm{Ri}_c,\ \mathrm{Ri}_c^\ast$
 & Threshold for turbulence collapse. $\mathrm{Ri}_c^\ast$ denotes dynamic, state-dependent value (symbolic regression target). \\

Neutral curvature invariant
 & $\Delta$
 & $\Delta = a_h - 2 a_m$. Determines the sign of $\partial_z^2 \mathrm{Ri}_g$ near neutrality, which sets the initial sign of the bias ratio $B$. \\

Turbulent kinetic energy
 & TKE
 & $0.5(\overline{u'^2}+\overline{v'^2}+\overline{w'^2})$. Used in higher-order closures for turbulence memory and intermittency. \\

\hline
\end{tabularx}
\caption{Canonical definitions and symbols used in the ABL Closure Project and associated machine-learning components.}
\label{tab:abl-glossary}
\end{table}